\subsection{Matematika v LaTeXu}

LaTeX je še posebej primeren za pisanje matematičnih izrazov in ena izmed njegovih glavnih prednosti je prav podpora za matematične simbole in formule.
LaTeX je že sam po sebi zelo zmogljiv. Vendar pa obstaja tudi veliko paketov, ki še dodatno razširijo njegove zmogljivosti na področju matematike.

Standardni paket za matematiko so:

\begin{itemize}
\item \texttt{amsmath}: paket, ki ga je razvila Ameriška matematična družba (American Mathematical Society) in ponuja številne izboljšave za pisanje matematičnih izrazov.
\item \texttt{amssymb}: paket, ki ponuja dodatne matematične simbole.
\item \texttt{amsfonts}: paket, ki omogoča uporabo dodatnih matematičnih pisav.
\item \texttt{amsthm}: paket za definiranje in oblikovanje matematičnih izrekov, trditve, definicij itd.
\end{itemize}

Ti paketi so pogosto vključeni v predloge dokumentov, vendar jih je mogoče vključiti tudi ročno z ukazom \texttt{{\textbackslash}usepackage\{amsmath, amssymb, amsfonts, amsthm\}} v preambuli vašega LaTeX dokumenta.

V sodoben LaTeX dokument je priporočljivo vključiti paket \texttt{mathtools}, ki je razširitev paketa \texttt{amsmath} in ponuja dodatne funkcije in izboljšave za pisanje matematičnih izrazov.

\begin{framed}
\textbf{Pomembno!}\\
V tem poglavju bomo predpostavljali, da so ti paketi že vključeni v vaš dokument. Lahko jih vključite z naslednjim ukazom v preambuli vašega LaTeX dokumenta:

\begin{verbatim}
\usepackage{mathtools}  % mathtools vključuje tudi amsmath
\usepackage{amssymb, amsfonts, amsthm}
\end{verbatim}
\end{framed}

\subsubsection{Osnovni matematični izrazi}

LaTeX omogoča pisanje matematičnih izrazov na dva načina: znotraj besedila (\textit{vrstični način}; ang. \textit{inline mode} ) in v ločenih vrsticah (\textit{prikazen način}; ang. \textit{display mode}).

Matematika v vrstičnem načinu se vstavlja z uporabo znakov \texttt{\$...\$} ali \texttt{{\textbackslash}(...{\textbackslash})}. Na primer, izraz za Pitagorov izrek lahko zapišemo kot \texttt{\$a\^2 + b\^2 = c\^2\$} (ali \texttt{{\textbackslash}(a\^2 + b\^2 = c\^2{\textbackslash})}), ki je natisnjen kot $a^2 + b^2 = c^2$.

Za prikaz matematičnih izrazov v ločenih vrsticah (in osredotočeno poravnano) se uporabljajo znaki \texttt{{\textbackslash}[...{\textbackslash}]} ali okolje \texttt{equation*} (sta enakomerna). Na primer, če je $ax^2 + bx +c = c$ kvadratna enačba, potem je rešitev te enačbe:

\begin{verbatim}
\[x_{1,2} = \frac{ -b \pm \sqrt{b^2 - 4ac}}{2} \]
\end{verbatim}

kar je natisnjeno kot

\begin{equation}
x_{1,2} = \frac{-b \pm \sqrt{b^2 - 4ac}}{2}.
\end{equation}

Ta izraz v vrstičnem načinu bi bil zapisan kot \texttt{\$x\_\{1,2\} = {\textbackslash}frac\{-b {\textbackslash}pm {\textbackslash}sqrt\{b\^2 - 4ac\}\}\{2\}\$}, kar je natisnjeno kot $x_{1,2} = \frac{ -b \pm \sqrt{b^2 - 4ac}}{2}$.

\begin{framed}
\textbf{Opomba}\\
Matematika v vrstičnem načinu zapisana z \texttt{\$...\$} včasih se šteje za zastarelo. Priporočljivo je uporabljati \texttt{{\textbackslash}(...{\textbackslash})}. Vendar, v praksi večinoma avtorjev uporabijo \texttt{\$...\$}, ker je hitreje za tipkanje in dobi se enak rezultat.

Podobno je tudi pri prikaznem načinu, kjer je priporočljivo uporabljati \texttt{{\textbackslash}[...{\textbackslash}]} (ali okolja \texttt{equation*}) namesto \texttt{\$\$ ... \$\$}s. Vendar je \texttt{\$\$...\$\$} resnično zastarelo in \textbf{se ne sme uporabljati}.
\end{framed}

Matematične izraze v prikaznem načinu je mogoče tudi oštevilčiti z uporabo okolja \texttt{equation} (brez zvezdice). Rezultat je podoben kot zgoraj, vendar bo enačba oštevilčena. Seveda, lahko sklicujemo na to enačbo v besedilu z uporabo ukazov \texttt{{\textbackslash}label\{\textless oznaka\textgreater \}} in \texttt{{\textbackslash}ref\{\textless oznaka\textgreater \}}:

\begin{example}\label{eg_matematika_equation_1}\end{example}Za enačbe lahko tudi uporabimo ukaze \texttt{{\textbackslash}eqref\{\textless oznaka\textgreater \}} in \texttt{{\textbackslash}tag\{\textless opis\textgreater \}}, ki omogočata bolj prilagojeno sklicevanje in označevanje enačb:

\begin{example}\label{eg_matematika_equation_2}\end{example}\begin{example}\label{eg_matematika_equation_3}\end{example}Ukazi in okolja za pisanje matematičnih aktivirajo matematični način, kjer so nekateri znaki in ukazi drugačni kot v običajnem besedilnem načinu. Na primer, presledki v matematičnem načinu nimajo pomena, medtem ko so v besedilnem načinu pomembni:

Vse črke so v matematičnem načinu obravnavane kot spremenljivke in so zato natisnjene v poševni pisavi (italics). Razmik med črkami je večji kot v običajnem besedilu.

Zaradi zgoraj navedenih razlogov moramo matematično besedilo \textbf{vedno} zapisovati v matematičnem načinu (v vrstici ali prikazano).

Če želimo v matematičnem načinu uporabiti običajno besedilo, lahko uporabimo ukaz \texttt{{\textbackslash}text\{...\}}: why?

\subsubsection{Stavljenje matematičnih formul}

Možnosti za matematične simbole je ogromno, zato bomo prikazali le nekaj
primerov.
Seznam primerov tukaj je namenjen prikazu nekaterih najpogostejših primerov uporabe, vendar nikakor ni izčrpen.

Pri sestavljanju matematičnih formul vam bodo morda v pomoč naslednji viri.

\begin{itemize}
\item \href{https://tug.ctan.org/info/symbols/comprehensive/symbols-a4.pdf}{Comprehensive LaTeX Symbol List}: Obsežen seznam LaTeX simbolov, ki vključuje simbole iz različnih paketov \textbf{Skoraj 500 strani}.
\item \href{http://detexify.kirelabs.org/classify.html}{Detexify}: Orodje, ki vam omogoča, da narišete simbol, ki ga iščete, in orodje vam bo predlagalo ustrezne LaTeX ukaze za ta simbol.
\item \href{./img/04\_seznamSimbolov.pdf}{Seznam osnovnih matematičnih simbolov}\{target=``\_blank''\}: Kratek seznam pogosto uporabljenih LaTeX matematičnih simbolov.
\end{itemize}

\subsubsection{Izreki, dokazi, definicije, itd..}

Ko sestavljanju matematičnih besedil je pogosto potrebno vključiti izreke, dokaze, definicije in podobne strukture (ki jih bomo zaradi enostavnosti imenovali ``teoremi'').
Paket \texttt{amsthm}, ponuja enostaven način za ustvarjanje teh struktur z uporabo okolij.

Najprej moramo definirati nova okolja za naše teoreme v preambuli dokumenta.
To storimo z uporabo ukaza

\begin{verbatim}
\newtheorem{<ime_okolja>}[<drugi_teorem>]{<Napis>}[<razdelek>]
\end{verbatim}

kjer je \texttt{\textless ime\_okolja\textgreater } ime novega okolja in
\texttt{\textless Napis\textgreater } je besedilo, ki se bo prikazalo pred številko teorema.
Neobvezna argumenta \texttt{\textless drugi\_teorem\textgreater } in \texttt{\textless razdelek\textgreater } določata, kako se bodo teoremi številčili (in lahko uporabimo samo ena od njih ali nobene).
Uporabimo \texttt{\textless drugi\_teorem\textgreater }, če želimo, da se novi teoremi številčijo skupaj z obstoječim teoremom, kjer je \texttt{\textless drugi\_teorem\textgreater } ime okolja.
Če želimo, da se številčenje teoremov ponastavi v vsakem razdelku, uporabimo \texttt{\textless razdelek\textgreater }, kjer je \texttt{\textless razdelek\textgreater } ime števca razdelkov (npr. \texttt{section} ali \texttt{subsection}).
Če ne uporabimo nobenega od teh dveh argumentov, se bodo teoremi številčili neodvisno skozi celoten dokument (1, 2, 3, ...).
Ukaz \texttt{{\textbackslash}newtheorem*\{\textless ime\_okolja\textgreater \}\{\textless Naslov\textgreater \}} definira različico okolja, ki ne bo imela številke.

Po definiciji okolja lahko uporabimo novo okolje v našem dokumentu z uporabo standardnega LaTeX okolja sintakse:

\begin{verbatim}
\begin{<ime_okolja>}[<Naslov>]
Vsebina teorema.
\end{<ime_okolja>}
\end{verbatim}

Neobvezni argument \texttt{\textless Naslov\textgreater } omogoča dodajanje naslova teorema, ki bo prikazan poleg številke teorema.

\begin{example}\label{eg_theorems_1}\end{example}V LaTeXu obstajajo tri različne sloge teoremov:

\begin{itemize}
\item \texttt{plain}: privzeti slog, krepko naslov in poševno telo.
\item \texttt{definition}: krepko naslov in pokončno telo.
\item \texttt{remark}: poševno naslov in pokončno telo.
\end{itemize}

Za nastavitev sloga teorema uporabimo ukaz

\begin{verbatim}
\theoremstyle{<slog>}
\end{verbatim}

pred definicijo okolja teorema.
Vsak teorem, ki je definiran po tem ukazu, bo uporabil določen slog.

\begin{example}\label{eg_theorems_2}\end{example}Paket \texttt{amsthm} ponuja tudi okolje \texttt{proof} za pisanje dokazov.
Okolje samodejno doda besedo ``Dokaz'' na začetek in simbol kvadrata na konec dokaza.

\begin{example}\label{eg-theorems_3}\end{example}Ukaz \texttt{{\textbackslash}qedhere} lahko uporabimo znotraj okolja \texttt{proof}, da postavimo simbol kvadrata na določeno mesto, običajno na konec enačbe.

\begin{example}\label{eg-theorems_4}\end{example}\subsubsection{Vaje}\label{latexMatematika_vaje}

\begin{itemize}
\item Exercise~\ref{ex-matematika-1}: Napišite LaTeX kodo za osnovne matematične izraze.
\item Exercise~\ref{ex-matematika-2}: Napišite LaTeX kodo za matematične izraze z uporaba okolji \texttt{align} in \texttt{align*} in \texttt{aligned}.
\item Exercise~\ref{ex-matematika-3}: Napišite LaTeX kodo, da reproducira Maxwellove enačbe.
\item Exercise~\ref{ex-matematika-4}: Napišite LaTeX kodo za izreke in dokaze.
\end{itemize}